\documentclass[11pt,a4paper]{article}
\usepackage[T1]{fontenc}
\usepackage[utf8]{inputenc}
\usepackage{lmodern}
\usepackage{amsmath,amssymb,bm}
\usepackage{geometry}
\geometry{margin=25mm}
\usepackage{booktabs}
\usepackage{array}
\usepackage{pdflscape}
\usepackage{caption}
\usepackage[hidelinks]{hyperref}
\newcommand{\CE}{\ensuremath{\mathrm{CE}}}
\setlength{\parindent}{0pt}
\setlength{\parskip}{0.55em}
\begin{document}
\begin{center}
{\LARGE\bfseries
Numerical Convergence of High-Order Harmonic Generation Simulations for Bayesian Optimization\par}

\vspace{1.1em}

{\normalsize
Dominika Maslarova$^{1,*}$, Mattys Pouyez$^{1}$,
Patrik Jansson$^{2}$, Jan Vabek$^{3,4}$,\\
Robin Weissenbilder$^{5}$, Linus Hagedorn$^{1,3}$,
Anne L'Huillier$^{5}$, and Tünde Fülöp$^{1}$\par}

\vspace{1em}

\begin{minipage}{0.95\textwidth}
\centering
\small

$^{1}$Department of Physics and Astronomy,
Chalmers University of Technology, Gothenburg, Sweden\par
\vspace{0.25em}

$^{2}$Department of Computer Science and Engineering,
Chalmers University of Technology and University of Gothenburg,
Gothenburg, Sweden\par
\vspace{0.25em}

$^{3}$Czech Technical University in Prague,
Faculty of Nuclear Sciences and Physical Engineering,
Prague, Czech Republic\par
\vspace{0.25em}

$^{4}$ELI Beamlines Facility,
The Extreme Light Infrastructure ERIC,
Doln\'{\i} B\v{r}e\v{z}any, Czech Republic\par
\vspace{0.25em}

$^{5}$Department of Physics,
Lund University, Lund, Sweden\par

\vspace{0.7em}

$^{*}$Corresponding author:
\texttt{dommas@chalmers.se}

\end{minipage}
\end{center}

\vspace{1.2em}


\section{Numerical convergence}\label{sec:numconv}

To assess the numerical accuracy of the high-order harmonic generation simulations performed with the HarP code, we investigated the sensitivity of the conversion efficiency (CE) to numerical discretization and computational-domain size. For the two-dimensional (2D) optimization, we tested 32 points distributed across the pressure--length domain. For the four-dimensional (4D) optimization, we examined the predicted optimum and selected additional configurations covering different laser parameters.

Preliminary tests indicated the greatest sensitivity to radial discretization and the smallest sensitivity to longitudinal discretization. Consequently, the discretization step sizes selected for the main calculations were
\begin{equation}
(\Delta r,\Delta z,\Delta t)=(202.9~\mathrm{nm},L/256,128.9~\mathrm{as}).
\end{equation}
In the numerical tests, these were compared with a higher-resolution configuration defined by
\begin{equation}
(\Delta r_{\mathrm{res}},\Delta z_{\mathrm{res}},\Delta t_{\mathrm{res}})
=(\Delta r/2,\Delta z/8,\Delta t/4).
\end{equation}

We also examined the influence of computational-domain size by doubling the radial and temporal extents from $L_r=3w_0$ and $L_t=3\tau$ to $L_r=6w_0$ and $L_t=6\tau$, while maintaining the original discretization steps.

The CEs obtained with the higher-resolution and enlarged-domain configurations are denoted by $\CE_{\mathrm{res}}$ and $\CE_{\mathrm{domain}}$, respectively. Their absolute differences relative to the CE obtained with the numerical configuration used in the optimization, $\CE_{\mathrm{final}}$, are defined as
\begin{equation}
\Delta\CE_{\mathrm{res}}=|\CE_{\mathrm{final}}-\CE_{\mathrm{res}}|,
\qquad
\Delta\CE_{\mathrm{domain}}=|\CE_{\mathrm{final}}-\CE_{\mathrm{domain}}|.
\end{equation}

In the 2D optimization, the 32 test points were generated using a Sobol sequence across the same pressure--length domain as in the accompanying study, with the laser intensity and pulse duration kept fixed. The results are summarized in Table~\ref{tab:numerical_convergence}. The maximum absolute differences were $4.57\times10^{-8}$ for the higher-resolution calculations and $1.30\times10^{-8}$ for the enlarged-domain calculations, both small compared with the characteristic CE values of order $10^{-6}$ in the high-efficiency regions relevant to the optimization.

Although larger relative differences were observed in regions of low CE, their absolute magnitudes remained small. For example, at point 28, increasing the numerical resolution resulted in a relative difference of approximately $13\%$, but an absolute difference of only $1.7\times10^{-9}$. Similarly, enlarging the computational domain produced a relative difference of approximately $50\%$, corresponding to an absolute difference of only $5.6\times10^{-9}$. These results show that even relatively large percentage differences in low-efficiency regions correspond to small absolute changes compared with the CE values relevant to the optimization.

In the 4D optimization, the same numerical configurations were tested at the Gaussian-process-predicted optimum and at additional points representing different regions of the parameter space. At the predicted optimum, increasing the numerical resolution changed CE by $9.0\times10^{-8}$, which corresponds to approximately $1.3\%$. Enlarging the computational domain produced a smaller absolute difference of $0.8\times10^{-9}$.

The additional 4D test points revealed a greater sensitivity to numerical resolution, particularly at high driving intensities and longer pulse durations. The largest observed absolute difference was approximately $1.2\times10^{-7}$ ($4.8\%$), at a driving intensity of $10^{15}~\mathrm{W/cm^2}$ and a pulse duration of $11~\mathrm{fs}$. The influence of computational-domain size remained considerably smaller, with relative differences below $0.12\%$ at the tested points.

To further investigate the numerical sensitivity, we refined each discretization step separately, reducing $\Delta t$, $\Delta r$, and $\Delta z$ by factors of four, two, and eight, respectively. At the best 2D point, these refinements changed CE by approximately $+0.6\%$, $-0.6\%$, and $+0.25\%$, respectively. At the 4D optimum, the corresponding changes were $+1.6\%$, $-1.2\%$, and $+0.7\%$. These results show a greater sensitivity to numerical resolution in 4D, with the effects of individual refinements partially compensating when combined.

Overall, the selected numerical parameters provide sufficiently stable results for the optimization studies, with relatively small differences in the high-efficiency regions of interest. Furthermore, the two best 4D configurations retained the same ranking at both numerical resolutions, supporting the stability of the identified optimum.



\begin{table}[p]
\centering
\caption{\textbf{Numerical convergence study at 32 Sobol points distributed across the pressure--length domain.} The normalized coordinates $L/z_R$ and $pz_R$ represent the medium length and gas pressure, respectively, where $L$ is the medium length, $p$ is the gas pressure, and $z_R$ is the Rayleigh length. Conversion efficiencies are reported for the selected final numerical configuration ($\CE_{\mathrm{final}}$), higher numerical resolution ($\CE_{\mathrm{res}}$), and enlarged computational domain ($\CE_{\mathrm{domain}}$). Absolute differences quantify sensitivity to numerical discretization and computational-domain size. All CE values and differences are given in units of $10^{-8}$.}
\label{tab:numerical_convergence}
\vspace{0.8em}
\normalsize
\renewcommand{\arraystretch}{1.00}
\setlength{\tabcolsep}{6pt}
\begin{tabular}{r r r r r r r r}
\toprule
Point & $L/z_R$ & $pz_R$ [mbar cm] & $\CE_{\mathrm{final}}$ & $\CE_{\mathrm{res}}$ & $\CE_{\mathrm{domain}}$ & $\Delta\CE_{\mathrm{res}}$ & $\Delta\CE_{\mathrm{domain}}$ \\
\midrule
1 & 1.00 & 14.86 & 101.98 & 102.16 & 101.69 & 0.19 & 0.29 \\
2 & 0.65 & 296.12 & 3.46 & 3.37 & 3.47 & 0.09 & 0.01 \\
3 & 0.12 & 126.70 & 273.84 & 272.40 & 273.85 & 1.44 & 0.01 \\
4 & 1.26 & 183.99 & 1.45 & 1.48 & 1.44 & 0.02 & 0.01 \\
5 & 1.44 & 88.15 & 2.46 & 2.40 & 2.31 & 0.06 & 0.14 \\
6 & 0.30 & 225.73 & 22.50 & 22.63 & 22.45 & 0.13 & 0.05 \\
7 & 0.47 & 51.45 & 39.00 & 38.39 & 38.92 & 0.61 & 0.09 \\
8 & 0.82 & 262.14 & 2.43 & 2.38 & 2.42 & 0.05 & 0.01 \\
9 & 0.87 & 142.24 & 5.17 & 5.09 & 5.15 & 0.08 & 0.01 \\
10 & 0.53 & 168.46 & 8.81 & 8.76 & 8.78 & 0.05 & 0.03 \\
11 & 0.25 & 31.60 & 169.29 & 168.93 & 169.29 & 0.36 & 0.00 \\
12 & 1.39 & 279.39 & 0.62 & 0.66 & 0.51 & 0.04 & 0.10 \\
13 & 1.22 & 71.88 & 5.57 & 5.56 & 5.45 & 0.00 & 0.12 \\
14 & 0.08 & 241.71 & 328.62 & 329.36 & 328.62 & 0.74 & 0.00 \\
15 & 0.69 & 109.70 & 8.47 & 8.46 & 8.45 & 0.01 & 0.02 \\
16 & 1.04 & 204.17 & 2.01 & 2.02 & 2.01 & 0.01 & 0.00 \\
17 & 1.12 & 99.05 & 4.44 & 4.43 & 4.43 & 0.00 & 0.00 \\
18 & 0.73 & 193.65 & 5.02 & 4.96 & 5.02 & 0.07 & 0.01 \\
19 & 0.05 & 61.50 & 136.56 & 135.99 & 136.56 & 0.58 & 0.00 \\
20 & 1.14 & 231.49 & 1.47 & 1.49 & 1.44 & 0.01 & 0.03 \\
21 & 1.33 & 23.41 & 55.01 & 54.72 & 53.71 & 0.29 & 1.30 \\
22 & 0.24 & 271.04 & 26.14 & 26.08 & 26.13 & 0.06 & 0.01 \\
23 & 0.53 & 134.32 & 10.43 & 10.42 & 10.41 & 0.01 & 0.01 \\
24 & 0.93 & 160.42 & 3.70 & 3.64 & 3.72 & 0.06 & 0.02 \\
25 & 0.78 & 41.22 & 23.51 & 23.46 & 23.43 & 0.05 & 0.08 \\
26 & 0.39 & 251.78 & 11.41 & 11.38 & 11.41 & 0.03 & 0.00 \\
27 & 0.38 & 77.65 & 29.45 & 29.18 & 29.52 & 0.27 & 0.07 \\
28 & 1.48 & 215.06 & 1.11 & 1.28 & 0.55 & 0.17 & 0.56 \\
29 & 1.27 & 118.64 & 2.32 & 2.31 & 2.33 & 0.01 & 0.01 \\
30 & 0.18 & 176.10 & 168.57 & 164.00 & 168.55 & 4.57 & 0.03 \\
31 & 0.59 & 6.53 & 24.23 & 24.15 & 24.28 & 0.09 & 0.05 \\
32 & 0.98 & 287.92 & 1.50 & 1.51 & 1.53 & 0.00 & 0.02 \\
\bottomrule
\end{tabular}
\end{table}

\end{document}
